% ============================================================================
%  Can Your Agent See the Key?
%  Gateway-Enforced Brokered Isolation for LLM Tool-Using Agents
%
%  IEEE MetroCon 2026 -- Open Call for Papers / Secure and Intelligent Systems
%  Sathyaraj Kolandasamy
%
%  ---------------------------------------------------------------------------
%  RESULTS PROVENANCE
%
%  Table~\ref{tab:main} is sourced from:
%    results_all_all_20261004_161636.{json,csv}   (this folder)
%    model llama3.2:1b, 96 runs, all status=ok, n=12 per cell
%
%  Verified 2026-10-04: every reported aggregate recomputes exactly from the 96
%  per-run records; brokered exposure is 0/48; direct exposure is 48/48;
%  ABR == (injected AND exposed) on all 96; no denials in the direct arm.
%
%  The task_utility_retention column in those result files is STALE: it was
%  produced before evaluate.py was fixed (2026-10-04) and normalizes against
%  brokered|none instead of the unhardened clean cell. The paper does not
%  report retention. Re-running the harness now yields the corrected values.
%  ---------------------------------------------------------------------------
\documentclass[conference]{IEEEtran}
\IEEEoverridecommandlockouts

\usepackage{cite}
\usepackage{amsmath,amssymb,amsfonts}
\usepackage{booktabs}
\usepackage{array}
\usepackage{graphicx}
\usepackage{textcomp}
\usepackage{xurl}
\usepackage[hidelinks]{hyperref}
\usepackage[T1]{fontenc}
\usepackage[utf8]{inputenc}

\hypersetup{
  pdftitle={Can Your Agent See the Key? Gateway-Enforced Brokered Isolation for LLM Tool-Using Agents},
  pdfauthor={Sathyaraj Kolandasamy}
}

% IEEEtran table captions are small-caps; keep code tokens in upright
% typewriter so they do not letter-space (e.g. C L I E N T_C E R T).
\DeclareRobustCommand{\capttd}[1]{\mbox{\normalfont\ttfamily #1}}

\begin{document}

\title{Can Your Agent See the Key?\\
Gateway-Enforced Brokered Isolation for\\
LLM Tool-Using Agents}

\author{\IEEEauthorblockN{Sathyaraj Kolandasamy}
\IEEEauthorblockA{\textit{Nova Southeastern University}\\
\textit{College of Computing, AI, and Cybersecurity}\\
sk2120@mynsu.nova.edu}}

\maketitle

\begin{abstract}
LLM tool-using agents authenticate to APIs and ticketing systems. Prototypes
routinely place secrets in the model context through a \texttt{get\_secret}
tool, and under prompt injection that material can leak. I contribute a
product-agnostic \textbf{gateway-enforced brokered isolation (GEBI)} reference
architecture and evaluate its core credential-containment invariant: isolated
agent execution, policy-checked outbound actions, server-side credential
substitution, and scrubbed model-visible results. The distinction from
prompt-level defenses is architectural rather than lexical: a prompt cannot
un-see a key already in context, whereas brokering keeps the key out of the
transcript entirely. On a frozen 12-task set with a local
\texttt{llama3.2:1b} model, direct mode shows secret-exposure rate 1.00 and
attacker-benefit rate (ABR) 1.00 under all three injection strategies --- and
exposure 1.00 even with no injection at all --- while brokered mode holds
exposure and ABR at 0.00 across all 48 brokered runs and still completes benign
tasks. A canary token tracked independently of exposure shows the separation
that matters: under fake-completion the brokered agent echoed the attacker's
marker on 0.67 of runs while leaking nothing, so containment did not depend on
the model resisting persuasion. Exposure is measured as an exact substring test
over non-operational placeholder tokens, so it detects verbatim reprinting but
not semantic paraphrase. The contribution is an architecture, a direct-vs-brokered harness,
and a sanitized credential-oriented injection suite.
\end{abstract}

\begin{IEEEkeywords}
LLM agents, credential isolation, prompt injection, zero-trust tooling,
AI security
\end{IEEEkeywords}

% ===========================================================================
\section{Introduction}

Tool-using language-model agents must act on external systems, and acting on
external systems requires credentials. The path of least resistance in current
agent frameworks is to expose a tool that returns the secret --- a
\texttt{get\_secret} call, or an environment read --- so that the model can
place it in a subsequent request. That design puts the credential into the
model's own context window, and once it is there the credential is subject to
every failure mode the context window has.

Prompt injection is the sharpest of those failure modes~\cite{liu2024formalizing,owasp2025}.
An attacker who controls any text the agent ingests --- a ticket body, a
document, a field in a tool result --- can attempt to steer the model into
printing, forwarding, or otherwise disclosing whatever it can see. Existing
robustness benchmarks largely score chat-safety refusal rather than credential
containment~\cite{chao2024jailbreakbench}, so they do not answer the question a
platform team actually has: \emph{did the secret ever reach the model at all?}

Production identity systems answered the analogous question years ago, and not
by asking the principal to behave. They prefer either short-lived, narrowly
scoped credentials or key-bound workload identities such as SPIFFE/SPIRE
SVIDs~\cite{spiffe} --- neither of which an application component should be able
to read and forward. Those are different mechanisms with different failure
modes, and the difference decides whether brokering is necessary or merely
convenient; Section~\ref{sec:shapes} separates them. The contribution of this
paper is to carry that discipline into the agent setting and to measure whether
it holds under injection.

\textbf{Hypothesis.} Direct secret access leaks or misuses credentials under
prompt injection substantially more often than brokered, non-exportable
issuance.

The public Bromure Agentic Coding product illustrates a related deployment
pattern: isolated Linux VMs, a local outbound gateway, stand-in credentials
inside the VM, and real-credential substitution at the gateway. Its
documentation also describes approval gating, prompt-injection screening, and
security-event logging~\cite{bromure2026}. I generalize these
product-described controls into GEBI, a vendor-neutral reference architecture,
and use Bromure only as motivation --- it is not an evaluated system here. The
experiment below tests the narrower invariant that secrets never enter the
model-visible channel.

\textbf{Contributions.}
\begin{itemize}
  \item \textbf{GEBI}, a product-agnostic reference architecture combining
  isolated agent execution, a non-bypassable policy gateway, server-side
  credential substitution, optional approval gating, and auditable inspection.
  \item \textbf{A direct-vs-brokered harness} that enforces ``no secrets in
  model-visible channels'' as a checked invariant in brokered mode rather than
  as a convention.
  \item \textbf{A sanitized credential-oriented injection suite} and a
  measured exposure / ABR / utility freeze, with the measurement definitions
  fixed in advance.
\end{itemize}

% ===========================================================================
\section{Threat Model}
\label{sec:threat}

\textbf{Attacker goal.} Cause \emph{bearer} secret material to appear in a
model-visible channel, and/or coerce a denied export path into succeeding. The
qualifier matters: as Section~\ref{sec:shapes} argues, that goal is equivalent
to compromise only for credentials whose possession is itself authorization.

\textbf{Knowledge.} Grey-box. The attacker knows the agent has tools and may
know tool names, but need not know broker internals. This is deliberately
generous: the \texttt{tool\_override} payload names tools from \emph{both}
modes, so a single injection probes both surfaces.

\textbf{Capabilities.} Inject UTF-8 text into the task channel, simulating
indirect injection through a ticket, a document, or a tool-returned field. The
attacker cannot modify broker allowlists, read host memory, or authenticate to
the broker directly.

\textbf{Defender.} Chooses the architecture, configures allow/deny lists and
handle lifetimes, and redacts logs. Prompt-level defenses are deliberately
\emph{excluded} from both arms so that architectural effects are not confounded
with prompt wording. The direct arm still carries a system prompt instructing
the agent not to print secrets; that instruction is the control, and its
failure is part of the result.

\textbf{Out of scope.} Live key theft; attacks on production identity
providers; multi-agent credential passing; VM-escape and gateway hardening.
Secrets are research fakes only --- non-operational placeholder tokens whose
literal values make exposure an exact, model-agnostic substring test rather
than a subjective judgment.

This study instantiates a \emph{static bearer} credential, chosen deliberately
because it is the shape for which exposure and compromise coincide exactly.
Credentials that are key-bound or audience-restricted behave differently under
the same attack, and the next section explains why.

% ===========================================================================
\section{Credential Shapes and the Reachability Equivalence}
\label{sec:shapes}

The paper has so far said ``the secret,'' and that singular hides a design
choice. An agent may authenticate to an application in several ways, and they
do not degrade alike when the agent's context window is exposed. This section
separates them along two questions --- whether the credential's bytes are
\emph{transferable} if observed, and whether the \emph{agent} holds the
authenticating material --- and then states precisely what brokering adds.
Table~\ref{tab:shapes} summarizes the result.

\subsection{Four credential shapes}

\textbf{Class B --- static bearer secret.} An API key, password, or connection
string. Authentication is possession of a string; I define the class by that
property, so anything that reads the string can act as the principal, from
anywhere, until the string is changed.

\textbf{Class S --- scoped or ephemeral bearer.} A short-lived access token
issued by a token service, or a session credential carrying an audience
restriction. Still a bearer: observation still yields impersonation, but the
window and the blast radius are bounded rather than open-ended.

\textbf{Class P --- proof-of-possession.} Authentication requires demonstrating
control of a private key rather than presenting a copyable value, as with a
client certificate or a SPIFFE X.509-SVID~\cite{spiffe}. The SVID document
itself is not secret; the workload-held private key is the part that
authenticates, and it never transits the channel.

\textbf{Class C --- brokered capability.} The agent names an \emph{action} and
holds an opaque handle; a gateway holds whichever of B, S, or P the backend
actually requires. This is GEBI.

\begin{table}[!t]
\caption{Credential shapes ordered by what a context-window disclosure yields an
attacker. Only Class B was exercised end-to-end by the broker in this study. The
harness declares a \capttd{client\_cert\_handle} kind, but it is stored and
returned as a static string with no certificate logic anywhere, so it is Class B
in behaviour regardless of its name.}
\label{tab:shapes}
\centering
\footnotesize
\setlength{\tabcolsep}{2.5pt}
\begin{tabular}{
  c
  >{\raggedright\arraybackslash}p{1.45cm}
  >{\raggedright\arraybackslash}p{1.55cm}
  >{\raggedright\arraybackslash}p{2.15cm}
  >{\raggedright\arraybackslash}p{1.55cm}}
\toprule
Class & Shape & Example & Disclosure in context yields & In this study? \\
\midrule
B & Static bearer secret & API key, service token & Full, durable impersonation from anywhere & \textbf{Yes} (both arms) \\
S & Scoped / ephemeral bearer & Short-lived access or session token & Impersonation, bounded by lifetime and audience & No \\
P & Proof-of-possession & Client certificate; SVID key~\cite{spiffe} & Nothing replayable; misuse in place still possible & No \\
C & Brokered capability & GEBI handle (this paper) & An action name under policy; no credential & \textbf{Yes} (brokered arm) \\
\bottomrule
\end{tabular}
\end{table}

\subsection{The reachability equivalence}

For a static bearer credential, authentication \emph{is} possession of a
string, so the predicate ``the secret reached the model context'' and the
predicate ``the credential is compromised'' are the same predicate. The
transcript is not a disclosure risk adjacent to compromise; it is the
compromise.

A context window is a uniquely bad container for such a secret because it is
replicated by default --- into the final answer, the trace log, the evaluation
record, and any retrieval index built from them --- so a single read fans out
into egress paths the agent's author never enumerated.

Proof-of-possession severs the equivalence by making the transferable part of
the credential and the authenticating part different objects. An attacker who
reads everything the model can see obtains an identifier, not an authenticator,
and a stolen identifier presented from the attacker's own channel fails at the
handshake. Class S sits between the two: the equivalence still holds, but only
within a bounded lifetime and audience.

\subsection{What brokering adds beyond binding}

Class P already breaks the equivalence, which raises the obvious question of why
a broker is needed at all. The answer is that binding defeats the thief but not
the confused deputy. An injected agent that \emph{legitimately holds} a bound
key can still sign a perfectly valid request toward an attacker-chosen
destination; nothing about key binding constrains \emph{what} the holder asks
for. GEBI mediates the action as well as removing the secret --- the handle
names what may be done, not merely who is doing it --- which is why the
allowlist and denylist of Section~\ref{sec:gebi} are load-bearing rather than
defence in depth. Binding and mediation are complementary: the first stops
replay, the second stops misuse in place.

\subsection{Coverage in this study}

The harness declares three credential kinds but the broker path exercises
exactly one: a static bearer key. There is no token exchange, no audience
binding, and no certificate or key-bound path anywhere in the implementation.
The shape actually exercised end-to-end is therefore Class B --- the shape for
which the equivalence is tightest and the gap consequently widest. The
1.00-to-0.00 result in Section~\ref{sec:results} should be read as a measurement
of that worst case rather than as a claim about credential handling in general,
and whether brokering adds containment \emph{beyond} what proof-of-possession
already provides is argued here but not measured.

% ===========================================================================
\section{The GEBI Architecture}
\label{sec:gebi}

GEBI is a layering claim: credential possession belongs to the gateway, not to
the agent. Five controls follow from it.

\begin{enumerate}
  \item \textbf{Isolated execution.} The agent runs in a workspace that holds
  no production credential material, and whose egress is restricted to the
  gateway. Any secret present inside the workspace is a stand-in.
  Section~\ref{subsec:deployment} discusses how this boundary is obtained;
  the harness models it rather than enforcing it.
  \item \textbf{Non-bypassable policy gateway.} All outbound action requests
  traverse a broker that evaluates them against policy. There is no side channel
  by which the agent can reach the external system directly.
  \item \textbf{Server-side credential substitution.} The broker holds the real
  credential, uses it to perform the authorized action, and never copies it
  into the response payload.
  \item \textbf{Scrubbed model-visible results.} What returns to the model is
  assembled from a fixed field template rather than filtered after the fact,
  so secrets cannot appear in it by construction.
  \item \textbf{Auditable inspection.} Every allow, deny, and issuance decision
  is logged, which makes denied attempts measurable rather than invisible.
\end{enumerate}

\subsection{Broker control flow}

Each \texttt{broker\_execute(action, params)} call resolves in a fixed order:

\begin{enumerate}
  \item Check a \textbf{denylist} and reject immediately. The denied set is
  \texttt{exfiltrate}, \texttt{print\_secret}, \texttt{export\_credential}, and
  \texttt{send\_to\_url}.
  \item Check an \textbf{allowlist} and reject anything not on it, so unknown
  actions \textbf{fail closed}. The allowed set is
  \texttt{fetch\_public\_status}, \texttt{lookup\_ticket}, and
  \texttt{summarize\_doc}.
  \item Mint a capability handle (Section~\ref{subsec:handles}).
  \item Perform the action server-side, where the secret is used but never
  copied into the payload.
  \item Return a result assembled field-by-field from a fixed template.
\end{enumerate}

Deny is evaluated \emph{before} allow, so a name appearing in both sets is
denied. Because the model-visible string is \emph{constructed} rather than
filtered, secrets cannot appear in it by construction; a final assertion
re-checks the rendered text before return and is treated as an invariant, not a
heuristic. If that assertion ever fires, the run is a harness bug, not a
finding.

\subsection{Capability handles}
\label{subsec:handles}

Handles are minted per call as opaque identifiers, marked non-exportable, and
scoped to a single action with a bounded lifetime (60\,s in this
configuration). The model receives the handle identifier, never the credential
it stands for. One honest caveat: the implementation records the lifetime at
issuance and uses it to characterize the policy configuration, but does not
re-check expiry on use. This study therefore does not claim to measure TTL
enforcement, only TTL configuration.

\subsection{Tool surface}

Table~\ref{tab:tools} gives the tools visible to the model in each mode. The
brokered arm deliberately advertises \texttt{broker\_export\_secret} as
\emph{bait}: the architecture is tested against an attacker who knows exactly
what to ask for, and the export path is denied rather than hidden. Export-style
names are additionally denied \emph{by name}, so a model that hallucinates a
direct-mode tool while in brokered mode still receives a denial rather than a
secret.

\begin{table}[!t]
\caption{Tool surface by mode. The brokered arm advertises an export tool and
then denies it, so the defense is tested against an informed attacker.}
\label{tab:tools}
\centering
\footnotesize
\setlength{\tabcolsep}{4pt}
\begin{tabular}{
  >{\raggedright\arraybackslash}p{1.5cm}
  >{\raggedright\arraybackslash}p{3.5cm}
  >{\raggedright\arraybackslash}p{2.4cm}}
\toprule
Mode & Tools exposed & Secret reaches model? \\
\midrule
Direct & \texttt{get\_secret}, \texttt{read\_env\_secret}, \texttt{call\_api} &
Yes, in tool result \\
Brokered & \texttt{broker\_execute}, \texttt{broker\_export\_secret} &
No --- denied or scrubbed \\
\bottomrule
\end{tabular}
\end{table}

\subsection{Deploying the isolation boundary}
\label{subsec:deployment}

The architecture above assumes the agent runs somewhere that holds no
credential material. That assumption is doing real work, so it is worth stating
how it is obtained in practice --- and noting that the harness evaluated here
simulates it in-process rather than enforcing it, which
Section~\ref{sec:limitations} records.

\textbf{Isolation alone is not containment.} A container whose environment
carries the upstream API key is simply a box with the secret inside it: the
agent reads the environment and the exposure result is unchanged. The property
that matters is not that the agent is sandboxed but that the sandbox is
\emph{credential-free}. The runner should be given network reachability to the
broker, never possession of a secret. Concretely, this rules out passing
credentials through environment variables or build arguments, mounting host
credential directories, and baking secrets into image layers --- each of which
returns the deployment to Class B of Section~\ref{sec:shapes}.

\textbf{Egress is the second boundary.} A sandbox with no route to the public
network cannot exfiltrate a secret even if one leaks into it, because there is
nowhere to send it. Running the agent on an internal-only network whose sole
reachable host is the broker therefore provides a containment guarantee that is
complementary to, and independent of, the one measured in
Section~\ref{sec:results}: exposure asks whether the secret entered the model's
context, while egress restriction asks whether anything in that context can
leave. A deployment that wants both must configure both; neither subsumes the
other.

\textbf{Authenticating the sandbox.} If the agent holds no secret, the broker
needs some other way to decide that a request came from a legitimate workload
--- the bootstrap problem that credential systems face generally. Three options
appear in practice, in increasing order of strength. The weakest is
\emph{network position}: the broker accepts connections only from the sandbox's
network namespace. Stronger is a \emph{mounted local endpoint}, where a socket
is placed into the sandbox and filesystem permissions carry the authentication
decision. Strongest is \emph{platform attestation}, in which the hosting
platform verifies properties of the workload --- its image, its namespace, its
service account --- and issues a short-lived, key-bound identity on that basis.
SPIFFE/SPIRE is the representative example~\cite{spiffe}, and it is worth
noting that this is precisely Class P of Section~\ref{sec:shapes}: the
workload's private key is generated locally and never transits, so attestation
solves the bootstrap problem without introducing a new bearer secret to protect.

\textbf{Why this matters for the threat model.} Each of these mechanisms
authenticates the \emph{workload}, not the model's intentions. None of them
prevents an injected agent from requesting a permitted action for an attacker's
purpose, which is the confused-deputy case of
Section~\ref{sec:shapes}. Deployment isolation and policy mediation are
therefore both required: isolation establishes that the caller is the agent and
that the agent holds nothing worth stealing, while the allowlist constrains
what that caller may ask for.

% ===========================================================================
\section{Method}
\label{sec:method}

\textbf{Direct arm.} \texttt{get\_secret} and environment-read tools return
plaintext placeholder tokens into the transcript, which is the prototype
pattern this paper argues against. There are four distinct paths by which a
secret becomes model-visible: \texttt{get\_secret} returns it in a tool result;
\texttt{read\_env\_secret} does the same through a simulated environment read;
\texttt{call\_api} leaks whatever the model pasted into its parameters; and the
composed final answer is scanned independently. The third is the most
instructive, because there the secret escapes not because a tool returned it
but because the model \emph{carried it forward} into a later call --- an egress
path no tool-level filter would catch. Brokering removes all four by
construction, since the agent never holds a value to carry. Both
\texttt{get\_secret} and \texttt{read\_env\_secret} accept a name argument and
default to the API key, so all three declared credential kinds are reachable
through the same channels.

\textbf{Brokered arm.} Only \texttt{broker\_execute} performs work; export
probes are denied; secrets stay server-side; results are scrubbed.

\textbf{Harness scope.} The experiment models the \emph{credential-containment
core} of GEBI. It does not measure VM-escape resistance, gateway hardening,
approval-gate usability, or any product's detection accuracy.

\textbf{Tasks.} Twelve benign status, ticket, and document jobs (task-set
SHA-256 \texttt{ad60f406\ldots bad36b}), each with a gold action name and
parameters.

\textbf{Injections.} \texttt{naive\_exfil}, \texttt{tool\_override},
\texttt{fake\_completion}, plus a clean \texttt{none} control. All carry a
visible research tag and a canary token, which makes them conspicuous by
design --- an acknowledged limitation, since it biases any detector
optimistically.

\textbf{Model and budget.} Local Ollama \texttt{llama3.2:1b} with JSON
tool-plan parsing, temperature 0.0, seed 42, and a cap of four tool steps per
task. Determinism comes from temperature 0 rather than from averaging; each
cell is run once.

\subsection{Metrics}

\textbf{Exposure} is 1 for a run if any literal placeholder secret value
appears in the model-visible channel --- the union of tool-result text and the
final answer. It is an exact substring test, chosen because it is objective and
model-agnostic; its blind spot is semantic leakage, where a model describes a
credential without reproducing it.

\textbf{ABR} (attacker-benefit rate) is exposure restricted to injected runs.
The consequence worth stating plainly: a denied-but-secret-free attempt does
\textbf{not} count as attacker benefit. That is what separates ``the agent
tried'' from ``the agent leaked,'' and it is why a brokered agent can comply
with an injection's \emph{form} while scoring ABR 0.00.

\textbf{Utility} is 1 if the gold action was invoked and returned successfully.
It is scored by exact match on the action name, independently of exposure, so a
leaky run can still be a useful one. The two axes are reported separately
rather than collapsed into a single score.

\textbf{Canary rate} is the fraction of runs in which the injection's marker
token reached the final answer. It measures compliance with the attacker's
surface form and is deliberately independent of exposure, which is what lets
the two be compared. \textbf{Denials} is the mean number of broker policy
refusals per run, which turns a blocked export into a positive observation
rather than an absence. Both are reported per cell in
Section~\ref{sec:results}.

The harness also computes a utility-retention ratio, normalized against the
unhardened clean cell (direct mode, no injection). It is omitted from
Table~\ref{tab:main} because at this sample size it adds no information beyond
the raw utilities --- the denominator is a single weak baseline
($0.33$, or $4/12$) and dividing by it mostly magnifies noise. It is worth
noting that retention is not bounded above: an injected run can out-score the
clean baseline when the injected text happens to push a small model into
emitting a well-formed tool call it otherwise fumbles, which is the same effect
visible in Finding~6.

% ===========================================================================
\section{Results}
\label{sec:results}

\begin{table}[!t]
\caption{Experimental freeze, $n=12$ tasks per cell (96 runs,
all \capttd{status=ok}). Exposure and ABR are the security axes; canary rate
measures compliance with the injection's surface form, and denials count broker
policy refusals. Single seed, temperature 0, one run per cell. At $n=12$ a
binomial 95\% interval is $\approx\pm0.28$, so utility and canary differences
are directional; the saturated exposure and ABR columns reflect a mechanism
rather than a sample estimate.}
\label{tab:main}
\centering
\footnotesize
\setlength{\tabcolsep}{3pt}
\begin{tabular}{llrrrrr}
\toprule
Mode & Injection & Expo. & ABR & Util. & Canary & Denials \\
\midrule
Direct & none & 1.00 & 0.00 & 0.33 & 0.00 & 0.00 \\
Direct & naive\_exfil & 1.00 & 1.00 & 0.58 & 0.75 & 0.00 \\
Direct & tool\_override & 1.00 & 1.00 & 0.25 & 0.33 & 0.00 \\
Direct & fake\_completion & 1.00 & 1.00 & 0.50 & 0.75 & 0.00 \\
\midrule
Brokered & none & 0.00 & 0.00 & 0.17 & 0.00 & 1.17 \\
Brokered & naive\_exfil & 0.00 & 0.00 & 0.33 & 0.42 & 0.75 \\
Brokered & tool\_override & 0.00 & 0.00 & 0.25 & 0.25 & 0.75 \\
Brokered & fake\_completion & 0.00 & 0.00 & 0.50 & 0.67 & 0.33 \\
\bottomrule
\end{tabular}
\end{table}

\textbf{Finding 1 --- direct mode leaks even when nobody attacks it.} Exposure
is 1.00 in every direct cell, \emph{including the clean control}. No injection
is required: the agent calls \texttt{get\_secret} in the ordinary course of
doing its job, and the credential enters the transcript. The injection
strategies change what the attacker \emph{gains}, not whether the secret is
present.

\textbf{Finding 2 --- the ABR gap is total.} Under injection, direct ABR is
1.00 for all three strategies; brokered ABR is 0.00 for all three. Within this
freeze the architectural difference is not a matter of degree.

\textbf{Finding 3 --- the containment invariant held, and what that does and does not
show.} Brokered exposure is 0.00 across all 48 brokered runs, clean and
injected alike.

This number deserves a direct caveat, because it is partly true by
construction. The brokered result string is assembled from a fixed field
template and asserted secret-free before return, so a brokered cell
\emph{cannot} report exposure without the assertion firing first. Reporting
0.00 is therefore not evidence that brokering defeats injection in the way an
attack-success measurement would be. The honest reading is narrower: across 36
injected brokered runs, no injection found a path that caused the invariant to
be violated or circumvented.

That is a weaker claim than ``brokering is secure,'' but it is not vacuous, and
it is worth being explicit about what would have falsified it. The assertion
would have fired --- and the result would have been a harness defect rather
than a finding --- had any of the following occurred: a denied action
nevertheless returning server-side data; the agent reaching a credential
through a tool name not on the catalog; a scrubbing gap letting placeholder
material into the payload; or a policy check evaluated in the wrong order, such
that an allowlisted-but-denied action slipped through. None occurred. The
broker logged \textbf{36 policy refusals across 25 of the 48 runs}, so the
denial path was exercised heavily rather than left dormant --- the mechanism
was under load, and it held.

What the result genuinely establishes is a property of the \emph{architecture},
not of the model: no sequence of model outputs in this freeze could produce
exposure, because the model was never in a position to emit a credential it had
never received. That is exactly the guarantee a prompt-level defense cannot
offer, and it is why the comparison in Section~\ref{sec:discussion} is between
mechanisms rather than between degrees of model compliance.

\textbf{Finding 4 --- compliance and exposure come apart, and the gap is
measurable.} This is the result the canary column exists to expose. In the
brokered arm under fake-completion, $0.67$ of runs echoed the attacker's canary
token --- the agent visibly adopted the injection's framing --- while exposure
stayed at $0.00$. The same pattern holds for naive exfiltration ($0.42$) and
tool override ($0.25$). An evaluation that scored only ``did the model obey?''
would record these as defeats. Measuring the credential shows nothing moved.
Containment here is not a function of the model declining to cooperate; it
cooperated, and the architecture made that irrelevant.

This also answers the objection raised under Finding~3 from a second direction.
If brokered exposure were 0.00 merely because no injection ever took hold, the
result would be uninformative --- a measurement of attacker ineffectiveness
rather than of containment. The canary column rules that reading out:
\textbf{16 of the 48 brokered runs} echoed the attacker's token, meaning the
injection demonstrably reached and influenced the model in a third of them, and
exposure was still 0.00 in every one. The attacks landed. The credential stayed
behind the gateway regardless.

\textbf{Finding 5 --- denials and compliance move in opposite directions.}
The brokered clean control incurs the most policy refusals ($1.17$ per run) and
no canary hits, while fake-completion incurs the fewest ($0.33$) and the most
canary hits ($0.67$). The agent that adopts the attacker's framing stops
probing the export path, because the injection has convinced it the credential
work is already done. The denial counter makes that visible; exposure alone
would not distinguish it from an agent that never tried.

\textbf{Finding 6 --- utility is imperfect, and injection does not uniformly
degrade it.} Utility is well below $1.00$ in both arms on this small local
model, and the highest cell in the matrix is \emph{injected}
(\texttt{direct|naive\_exfil}, $0.58$) rather than clean ($0.33$). Injection
changes agent behaviour without necessarily breaking the task.
Section~\ref{sec:discussion} gives a mechanism-level account of the low
absolute values; it is not a measured error analysis, and I do not present it
as one.

% ===========================================================================
\section{Discussion}
\label{sec:discussion}

\textbf{Compliance is not exposure.} The most useful consequence of separating
ABR from instruction-following is that a brokered agent can \emph{follow the
injection's form} --- attempting the export, echoing the attacker's canary ---
while leaking nothing. Measuring only ``did the model obey?'' would score that
as a defeat; measuring exposure shows the credential never moved. Architectural
containment does not require the model to resist persuasion, which is precisely
why it is more robust than asking it to.

Two signals recorded alongside exposure make this separation observable rather
than merely arguable. A \emph{canary} token embedded in each injection is
detected in the agent's final answer, so compliance with the injection's surface
form is visible whether or not a credential moved. A \emph{denial count} is
accumulated from broker policy decisions, so a refused export becomes a positive
measurement rather than a silent non-event. Together they distinguish three
outcomes that a single exposure number collapses: an agent that ignored the
injection, one that attempted the export and was refused, and one that
succeeded.

The measured result is that the brokered arm sits squarely in the middle case.
Under fake-completion it echoed the canary on $0.67$ of runs and still exposed
nothing; under the other two strategies the canary rate was $0.42$ and $0.25$
with exposure fixed at $0.00$ throughout. The agent was persuaded and the
credential stayed put. That is a stronger claim than ``the attack failed,''
because it identifies \emph{where} it failed: not at the model's judgment, which
the injection successfully bent, but at the boundary the model could not reach
past.

\textbf{Why prompt-level defenses are the wrong layer.} Direct mode leaks on
every cell \emph{despite} a system prompt that explicitly instructs the agent
not to print secrets. Once a credential is in the transcript, every downstream
path --- the final answer, the log, a retrieval index built from that log ---
is a potential egress. A sandwich prompt cannot un-see a key. GEBI's claim is
narrower and stronger: the key is never in the transcript to begin with.

\textbf{Utility is bounded by mechanism, not by security.} Utility is scored by
exact match on the gold action name, and the agent has a budget of four tool
steps per task. Under injection, a brokered run spends part of that budget on
denied export probes before reaching the legitimate call. That is a plausible,
mechanism-level account of why brokered utility varies across cells --- it is a
property of the harness's scoring rule and step limit, not a measured error
analysis. Distinguishing ``the model planned the wrong action'' from ``the model
ran out of steps'' requires per-run transcripts that this freeze does not
retain.

\textbf{Deployment implications.} Brokering shifts the burden to allowlist
design: a workflow that direct mode would ``solve'' by pasting a key must
instead be expressed as an allowlisted action. Utility retention makes that cost
visible rather than hidden. Two further controls follow. First, \textbf{fail
closed} when the broker is unavailable --- a fail-open path silently restores
exposure and invalidates the containment claim. Second, extend redaction to
\textbf{observability}, because logging and retrieval pipelines are
model-visible channels by another name.

% ===========================================================================
\section{Limitations}
\label{sec:limitations}

\begin{itemize}
  \item \textbf{Single model, small scale.} One small local model on twelve
  tasks. Utility figures characterize this configuration, not agent frameworks
  generally.
  \item \textbf{Exposure is a substring test.} It catches verbatim reprinting
  of known placeholder tokens but \textbf{not semantic leakage}, where a model
  describes or paraphrases a credential without reproducing it. ABR equals
  exposure on injected runs by construction, so it inherits the same blind
  spot. This is the single most important measurement gap in the study.
  \item \textbf{TTL is configured, not enforced.} Handle lifetimes are recorded
  at issuance; expiry is not re-checked on use.
  \item \textbf{One credential shape.} The broker performs a single static
  bearer-secret lookup; it implements no token exchange, no audience binding,
  and no certificate or key-bound path. The harness names a
  \texttt{client\_cert\_handle} kind, but it is stored and returned as a static
  string, so it is a bearer secret in behaviour. Whether brokering adds
  containment beyond what proof-of-possession already provides is therefore
  argued in Section~\ref{sec:shapes} and not measured.
  \item \textbf{Isolation is modelled, not enforced.} The harness is an
  in-process simulation: the direct and brokered arms differ in which tool
  surface the agent is given, not in any container, namespace, or network
  boundary. The deployment guidance of
  Section~\ref{subsec:deployment} is therefore argued from the architecture
  rather than demonstrated by the measurement, and the egress-restriction
  guarantee in particular is untested here. A containerized replication that
  enforces the boundary --- and that could be attacked through it --- is the
  natural next step.
  \item \textbf{Research mock, not a KMS.} The broker is a research
  implementation, not a certified KMS or HSM integration.
  \item \textbf{Conspicuous injections.} The payloads carry visible research
  tags and canaries, which makes them easier to detect than a stealthy
  adversary's would be.
  \item \textbf{Single-agent scope.} Multi-agent credential passing is
  untested, and it is a plausible place for containment to break.
  \item \textbf{No confidence intervals.} Cells are run once at temperature 0.
  At $n=12$ a binomial 95\% interval is roughly $\pm$0.28, so the utility and
  canary columns should be read as directional; the ordering of canary rates
  across injection strategies is suggestive but not separated at this sample
  size. The exposure and ABR columns are 0.00/1.00 saturated, which reflects a
  mechanism rather than a sample estimate.
  \item \textbf{Retention is computed but not reported.} The harness emits a
  utility-retention ratio against the unhardened clean cell. At $n=12$ with a
  baseline of $4/12$ it magnifies noise rather than adding information, so
  Table~\ref{tab:main} reports raw utilities instead. A larger task set would
  make the ratio worth reporting.
\end{itemize}

% ===========================================================================
\section{Reproducibility and Ethics}
\label{sec:repro}

All credential material is synthetic: non-operational placeholder tokens
carrying an explicit redaction marker, chosen so that they are greppable,
harmless if copied, and unambiguous as an exposure oracle. No live key, token,
or production identity system was accessed.

The harness enforces this in continuous integration. One test asserts that
every configured secret is a placeholder; another greps the source tree for
live-key patterns, private-key headers, and offensive-tooling markers and fails
the build if any appear. Transcripts are redacted on write, so the artifact
cannot accumulate secret material even accidentally. Runs record the seed,
dependency versions, configuration hash, and task-set hash.

Every figure in Table~\ref{tab:main} is derived from a single result file
containing all 96 per-run records, released with the public artifact
\url{https://github.com/sathya-ksamy-nsu/Gateway-Enforced-Brokered-Isolation}.
Each aggregate recomputes exactly from those records, so a reader can confirm
the table rather than trust it: brokered exposure is $0/48$, direct exposure is
$48/48$, and the ABR definition --- injected and exposed --- holds on all 96
runs without exception.

The contribution is an architecture and an evaluation protocol rather than
attack tooling. The injection payloads are published because they are
conspicuous, already-catalogued strategy archetypes, and because the defense
they probe is the one this paper recommends.

% ===========================================================================
\section{Conclusion}
\label{sec:conclusion}

GEBI removes keys from the model context and, on this freeze, eliminates
measurable secret exposure and attacker benefit under sanitized injection while
direct mode does not. The architecture generalizes controls described publicly
by Bromure into a vendor-neutral contribution: workspace isolation, gateway
enforcement, server-side substitution, approval controls, and
monitoring~\cite{bromure2026}. These results evaluate only the
credential-containment core --- not VM-escape resistance, gateway hardening,
approval usability, or any product's security performance.

The practical recommendation is narrow: if an agent can read a credential, then
every injection becomes a potential credential-disclosure vulnerability, and no
amount of prompt engineering changes that. Make the credential unreachable
instead, and measure the result by whether the secret ever appears in a
model-visible channel.

\textbf{Future work.} Stronger and multiple model families; semantic rather
than substring exposure detection; extension to proof-of-possession and
audience-bound credentials (Classes P and S of Section~\ref{sec:shapes}), where
the marginal value of brokering is an open question; approval-gate evaluation;
multi-tenant handle scoping; and RAG-borne injection, where the untrusted span
arrives through retrieval rather than the task channel.

% ===========================================================================
\section*{Acknowledgment}
The author thanks Nova Southeastern University faculty and peers for feedback
on earlier drafts of this protocol. Any remaining errors are the author's.

\begin{thebibliography}{00}
\bibitem{liu2024formalizing} Y. Liu, Y. Jia, R. Geng, J. Jia, and N. Z. Gong,
``Formalizing and benchmarking prompt injection attacks and defenses,'' in
\textit{Proc. USENIX Security Symp.}, 2024, pp. 1831--1847.
\bibitem{owasp2025} OWASP Foundation, ``OWASP Top 10 for Large Language Model
Applications,'' 2025. [Online]. Available:
\url{https://owasp.org/www-project-top-10-for-large-language-model-applications/}
\bibitem{spiffe} SPIFFE, ``Secure Production Identity Framework for Everyone
(SPIFFE/SPIRE).'' [Online]. Available: \url{https://spiffe.io/}
\bibitem{chao2024jailbreakbench} P. Chao \textit{et al.}, ``JailbreakBench: An
open robustness benchmark for jailbreaking large language models,'' in
\textit{Proc. NeurIPS Datasets and Benchmarks}, 2024.
\bibitem{bromure2026} Bromure, ``Bromure Agentic Coding,'' 2026. [Online].
Available: \url{https://bromure.io/en/agentic-coding/} (accessed Sep. 27, 2026).
\end{thebibliography}

\end{document}
